\documentclass[../../main.tex]{subfiles}
\begin{document}

\begin{flushright}
\footnotesize\textit{【自動文字起こし・要確認】}
\end{flushright}

[解] (1) $f(x) = \sin x$とおく。$f'(x) = \cos x, f''(x) = -\sin x$から、$0 \le x \le \pi/2$では$f(x)$は上に凸である。\\
ここで、左辺第2項は、$t = \pi - x$により、($\dfrac{dt}{dx} = -1$)
$$
\int_\beta^\alpha \sin(\pi-t) dt (-1)
$$
$$
= \int_\alpha^\beta \sin t dt
$$
となるから、$\sin(\pi-\beta) = \sin\beta$とあわせて、
$$
\int_0^\beta \sin t dt > \dfrac{1}{2}(\beta-\alpha)(\sin\alpha + \sin\beta) \cdots \text{\textcircled{1}}
$$
を示せば良いが、$0 \le \alpha < \beta \le \pi/2$及び右図の面積を比較して、\\
示すべきことは明らかである。

\paragraph{(2)}

(1)で $(\alpha, \beta) = \left( \dfrac{k\pi}{8}, \dfrac{k+1}{8}\pi \right) \ (k=0,1,2,3)$とすると、これは$0 \le \alpha < \beta \le \dfrac{\pi}{2}$をみたすので、
$$
\int_{\dfrac{k\pi}{8}}^{\dfrac{k+1}{8}\pi} \sin t dt > \dfrac{1}{2} \cdot \dfrac{\pi}{8} \left( \sin\dfrac{k\pi}{8} + \sin\dfrac{k+1}{8}\pi \right)
$$
$k$について足して、$\sin 0 = 0$から、
$$
\int_0^{\dfrac{\pi}{2}} \sin t dt > \dfrac{\pi}{16} \sum_{k=1}^7 \sin\dfrac{k\pi}{8} \quad (\because \sin(\pi-k) = \sin k)
$$
$$
\dfrac{16}{\pi} > \sum_{k=1}^7 \sin\dfrac{k\pi}{8}
$$

\end{document}
